\section{Metodologia e desenvolvimento da solução}
\label{sec:metodologia}

\subsection{Arquitetura do programa}
A solução proposta foi desenvolvida em Java, com uso das bibliotecas gráficas nativas da linguagem (Swing e AWT). A arquitetura é composta pelas classes \texttt{Main}, \texttt{JanelaPrincipal}, \texttt{Ponto}, \texttt{JanelaMundo}, \texttt{Transformacao} e \texttt{Display}, e pela enumeração \texttt{TipoNDC}.

A classe \texttt{Main} é responsável pela inicialização do programa. A \texttt{JanelaPrincipal} concentra a interface gráfica e coordena a interação com o usuário, a classe \texttt{Ponto} representa coordenadas, \texttt{JanelaMundo} armazena os limites definidos para o espaço de mundo, a classe \texttt{Transformacao} concentra as transformações utilizadas durante o programa, já a classe \texttt{Display} é responsável pela representação gráfica e pela ativação do pixel correspondente à posição calculada. A enumeração \texttt{TipoNDC} identifica os dois sistemas normalizados disponíveis, \texttt{ZERO\_A\_UM} (NDC \mbox{$[0,1] \times [0,1]$}) e \texttt{MENOS\_UM\_A\_UM} (NDC \mbox{$[-1,1] \times [-1,1]$}).

\subsection{Entradas do sistema}
A aplicação aceita duas formas de entrada, que percorrem o pipeline de transformações em sentidos diferentes.

\paragraph{Clique do mouse.} A primeira entrada é a posição de um clique do mouse sobre o display. O display possui resolução de $800 \times 600$ pixels, sendo que a origem das coordenadas começa no canto superior esquerdo.

Quando o usuário clica sobre o display, a aplicação obtém as coordenadas $x$ e $y$ fornecidas pelo evento de clique do mouse e cria um objeto do tipo \texttt{Ponto} com tais valores. A coordenada obtida, $P_{inp} = (x_{inp}, y_{inp})$, é utilizada como entrada para as transformações. O clique é capturado por um \textit{listener} associado ao display, de modo que cada clique substitui o pixel ativo anterior pela nova posição.

\paragraph{Consulta de ponto do mundo.} A segunda entrada é um ponto $P_w = (x_w, y_w)$ digitado pelo usuário no painel \textit{Consultar Ponto do Mundo}. Ao pressionar o botão \textit{Mostrar no display}, o ponto é convertido para coordenadas de dispositivo e o pixel correspondente é ativado. Essa entrada permite verificar diretamente o mapeamento Mundo $\rightarrow$ DC para valores escolhidos, como os do exemplo numérico da Seção~\ref{sec:fundamentacao-teorica}.

Em ambos os casos, a interface apresenta os valores correspondentes à coordenada de dispositivo, ao NDC \mbox{[0,1]}, ao NDC \mbox{[-1,1]} e à posição equivalente na janela de mundo.

\subsection{Seleção do sistema NDC}
\label{subsec:selecao-ndc}
Um seletor (\textit{combo box}) no painel \textit{Sistema NDC} permite escolher qual dos dois cenários de normalização é utilizado no pipeline: \texttt{ZERO\_A\_UM}, correspondente ao cenário (i), ou \texttt{MENOS\_UM\_A\_UM}, correspondente ao cenário (ii). O valor inicial é \texttt{ZERO\_A\_UM}.

A troca do sistema altera apenas quais funções de transformação serão utilizadas nas próximas interações; ela não dispara, por si só, uma nova ativação de pixel. Dessa forma, o pixel não é deslocado para um ponto que o usuário não solicitou, e o próximo clique ou consulta já utiliza o sistema escolhido. Independentemente da seleção, a interface sempre exibe o ponto nos dois sistemas NDC, o que permite compará-los.

\subsection{Transformações implementadas}
As transformações foram implementadas na classe \texttt{Transformacao}, aplicando diretamente as expressões já deduzidas na Seção~\ref{sec:fundamentacao-teorica} (Fundamentação teórica), que coincidem com as fórmulas da \aularef\ (inclusive por não inverterem o eixo $y$ entre o NDC e o dispositivo), para $W$ e $H$ correspondendo à largura e à altura do display. Quando o sistema selecionado é o NDC \mbox{$[0,1] \times [0,1]$}, o caminho utilizado para a ativação do pixel emprega as seguintes expressões:

\begin{itemize}
    \item \textbf{Entrada $\rightarrow$ NDC} (\texttt{inpToNdc01}), mesma forma da Equação~\eqref{eq:dc-ndc-i}, com $x_{inp}, y_{inp}$ no lugar de $x_d, y_d$:
    \[
        x_n^{(i)} = \frac{x_{inp}}{W-1}, \qquad y_n^{(i)} = \frac{y_{inp}}{H-1}.
    \]
    \item \textbf{NDC $\rightarrow$ Mundo} (\texttt{ndc01ToUser}), Equação~\eqref{eq:ndc-user-i}:
    \[
        x_w = x_{wmin} + x_n^{(i)} \, (x_{wmax} - x_{wmin}), \qquad
        y_w = y_{wmin} + y_n^{(i)} \, (y_{wmax} - y_{wmin}).
    \]
    \item \textbf{Mundo $\rightarrow$ NDC} (\texttt{userToNdc01}), Equação~\eqref{eq:user-ndc-i}:
    \[
        x_n^{(i)} = \frac{x_w - x_{wmin}}{x_{wmax} - x_{wmin}}, \qquad
        y_n^{(i)} = \frac{y_w - y_{wmin}}{y_{wmax} - y_{wmin}}.
    \]
    \item \textbf{NDC $\rightarrow$ DC} (\texttt{ndc01ToDc}), Equação~\eqref{eq:ndc-dc-i} sem o arredondamento final, aplicado apenas na etapa de ativação do pixel (Seção~\ref{subsec:ativacao-pixel}):
    \[
        x_d = x_n^{(i)} \, (W-1), \qquad y_d = y_n^{(i)} \, (H-1).
    \]
\end{itemize}

Quando o sistema selecionado é o NDC \mbox{$[-1,1] \times [-1,1]$}, as mesmas quatro etapas são realizadas pelos métodos \texttt{inpToNdc11}, \texttt{ndc11ToUser}, \texttt{userToNdc11} e \texttt{ndc11ToDc}, conforme as Equações~\eqref{eq:dc-ndc-ii}, \eqref{eq:ndc-user-ii}, \eqref{eq:user-ndc-ii} e \eqref{eq:ndc-dc-ii}. Como demonstrado na Subseção~\ref{subsec:sintese}, os dois cenários levam ao mesmo pixel final; a escolha altera apenas os valores intermediários exibidos e o conjunto de funções exercitado.

Vale observar que a implementação não inverte o eixo $y$ entre o NDC e o dispositivo, como nas fórmulas da aula, e que o raster do Java tem origem no canto superior esquerdo, com $y$ crescendo para baixo. Assim, diferentemente da convenção adotada na Seção~\ref{sec:fundamentacao-teorica} (origem no canto inferior esquerdo, $y$ para cima), no display um ponto com $y_w$ maior é ativado em uma linha mais baixa, e a imagem aparece espelhada verticalmente.

\subsection{Ativação do pixel}
\label{subsec:ativacao-pixel}
Após a realização das transformações, as coordenadas finais do dispositivo podem possuir valores fracionários. Como a posição de um pixel é representada por coordenadas inteiras, os valores obtidos são arredondados utilizando \texttt{Math.round()}. As coordenadas resultantes são enviadas ao método \texttt{drawPixel()} pertencente à classe \texttt{Display}.

A classe utiliza um \texttt{BufferedImage} com resolução de $800 \times 600$ pixels e fundo preto. O método \texttt{drawPixel()} verifica se a posição está dentro dos limites do display e, caso esteja, altera diretamente a cor daquele pixel por meio do método \texttt{setRGB()}. O pixel ativo é representado pela cor verde (\texttt{\#00FF00}). Caso a posição esteja fora dos limites, o display apenas é limpo e nenhum pixel é ativado.

Antes de ativar o novo pixel, o display é limpo. Dessa forma, somente a posição correspondente à coordenada atual permanece destacada.

\subsection{Interface para definição da janela de mundo}
A aplicação possui uma pequena interface que permite ao usuário definir os limites da janela de mundo. São disponibilizados campos para:
\begin{itemize}
    \item X mínimo ($x_{wmin}$);
    \item X máximo ($x_{wmax}$);
    \item Y mínimo ($y_{wmin}$);
    \item Y máximo ($y_{wmax}$).
\end{itemize}
Inicialmente, os campos são preenchidos com os valores $x_{wmin} = -100$, $x_{wmax} = 100$, $y_{wmin} = -100$ e $y_{wmax} = 100$.

O usuário pode alterar esses valores e pressionar o botão \textit{Aplicar}. A aplicação realiza a conversão dos valores informados para o tipo \texttt{double} e cria uma nova instância de \texttt{JanelaMundo} com os limites fornecidos. A nova janela passa a ser utilizada a partir da próxima interação (clique ou consulta).

A classe \texttt{JanelaMundo} verifica se os valores fornecidos formam uma janela válida, exigindo que o valor mínimo seja menor que o valor máximo em ambos os eixos. Caso sejam informados valores não numéricos ou limites inválidos, a aplicação apresenta uma mensagem de erro ao usuário e mantém a janela anterior.

\subsection{Fluxo geral da solução}
De forma resumida, o funcionamento da solução pode ser representado por dois fluxos. O primeiro corresponde ao clique do mouse:
\[
    \text{Mouse} \Rightarrow \text{INP} \Rightarrow \text{NDC} \Rightarrow \text{USER} \Rightarrow \text{NDC} \Rightarrow \text{DC} \Rightarrow \text{Pixel}
\]
O mouse fornece a coordenada inicial utilizada pelo nosso programa, as funções da classe \texttt{Transformacao} realizam as conversões, a janela de mundo fornece os limites para a transformação, e o \texttt{Display} utiliza a coordenada final para ativar diretamente o pixel correspondente. O percurso completo, de ida ao mundo e volta ao dispositivo, garante que o pixel ativado seja o resultado das transformações no sistema NDC selecionado, e não simplesmente a posição bruta do clique.

O segundo fluxo corresponde à consulta de um ponto do mundo:
\[
    \text{USER} \Rightarrow \text{NDC} \Rightarrow \text{DC} \Rightarrow \text{Pixel}
\]
Em ambos os fluxos, a etapa NDC utiliza o sistema escolhido no seletor (Subseção~\ref{subsec:selecao-ndc}).

A aplicação apresenta as coordenadas intermediárias na interface, permitindo acompanhar visualmente o resultado das transformações realizadas.
